\scanpage{88}
% L1–7
(vérification immédiate par la série). En appliquant au deuxième membre de $(8)$, où $t\frac{{\Theta}'_v(0\mid t)}{\Theta'_v(0\mid-\frac1t)}$ est laissé sous forme indéterminée, on trouve une relation sur $H(t)$ : $\frac{H'(t)}{H(t)}=-\frac1{2t}$, d'où suit que $H(t)$ est proportionnel à $\frac1{\sqrt t}$; le facteur de proportionnalité s'obtient en faisant dans

% L7
\[
\Theta(v\mid t)=e^{-\frac{i\pi v^2}{t}}\Theta\left(\frac vt\,\middle|\,-\frac1t\right)\frac a{\sqrt t},
\]

% L8–9
$t=i$ (d'où $t=-\frac1t$) et $v=\unsure{0|o}$, on trouve $a=\sqrt i$, d'où enfin

% L10, boxed
\[
\boxed{\Theta(v\mid t)=e^{-\frac{i\pi v^2}{t}}\Theta\left(\frac vt\,\middle|\,-\frac1t\right)i\sqrt{\frac it}}
\tag*{\textcircled{\scriptsize 10}}
\]

% L11–13
(où la détermination de $\sqrt{\frac it}$ est celle qui correspond à $-\frac{\pi}{2}\leq\operatorname{Arg}\frac it\leq+\frac{\pi}{2}$). Comparant avec \textcircled{\scriptsize 8}, on trouve donc

% L14
\[
\Theta'_v(0\mid t)=\frac{i}t\sqrt{\frac it}\,\Theta'_v(0\mid-\frac1t)
\tag*{\textcircled{\scriptsize 11}}
\]

% L15–16
Qu'on exprime généralement, faisant $t=iT$. (\unsure{où} $\unsure{\Re|\mathcal R}T>0$)

% L17
\[
\sum_1^\infty(-1)^{n-1}(2n-1)e^{-(n-\frac12)^2\pi T}
=\frac1{T^{3/2}}\sum_{\unsure{1}}^\infty(\unsure{-1})^{n-1}(2n-1)e^{-(n-\frac1{\unsure{2|3}})\frac{\pi}{\unsure{T|1}}}
\tag*{\textcircled{\scriptsize 12}}
\]

% L18
\[
\text{(car}\qquad\Theta'_v(\unsure{v|0}\mid t)=2\pi\sum_1^\infty(-1)^{n-1}(2n-1)e^{(n-\frac12)^2i\pi t}\cos(2n-1)\unsure{i\pi v|\pi v}
\tag*{\textcircled{\scriptsize \unsure{12|17} bis}}
\]

% L19, foot
\[
\text{d'où}\qquad\Theta'_v(0\mid iT)=2\pi\sum_1^\infty(\unsure{-1|1})^{n-1}(2n-1)e^{-(n-\frac12)^2\pi T}
\]

% lower left
% archive sheet number: 87
